% Source: https://artofproblemsolving.com/community/c3h2944325_san_jose_math_circle_monthly_contest_201718_7th8th_grade_contest_45
% Stable identity: aops-topic-2944325-q01
(a) The midpoint is $(7,7)$ and $AB$ has slope $-1$, so the bisector has slope $1$ through $(7,7)$: $y=x$.

(b) $2x^2=100$ gives $P=(5\sqrt2,5\sqrt2)$ and $Q=(-5\sqrt2,-5\sqrt2)$.
\begin{asy}
usepackage("amsmath");
size(3.5cm);
defaultpen(fontsize(10pt));
draw(circle((0,0),10));pair A=(8,6),B=(6,8),P=(5*sqrt(2),5*sqrt(2)),Q=-P;draw(A--B);draw(Q--P,dashed);
dot(A);dot(B);dot(P);dot(Q);dot((0,0));label("$A$",A,E);label("$B$",B,N);label("$P$",P,NE);label("$Q$",Q,SW);label("$O$",(0,0),SE);
\end{asy}

(c) The perpendicular bisector of a chord passes through the centre, so $PQ$ is a diameter: $PQ=\boxed{20}$.
